\chapter{اللقاء الحادي والعشرون: تابع تفسير سورة البقرة (آيات المطاعم والبر وبداية آيات القصاص)}

\section{مقدمة: بداية المجلد الثالث}
السلام عليكم ورحمة الله وبركاته. الحمد لله رب العالمين، وأشهد أن لا إله إلا الله وحده لا شريك له، وأشهد أن محمداً عبده ورسوله. اللهم صل على محمد وعلى آل محمد كما صليت على آل إبراهيم إنك حميد مجيد، اللهم بارك على محمد وعلى آل محمد كما باركت على آل إبراهيم إنك حميد مجيد.

افتتح الشيخ هذا اللقاء ببيان أننا دخلنا المجلد الثالث من تفسير الإمام \rawi{ابن جرير الطبري}، وبدأ من قول الله تعالى في سياق المطاعم والشرائع وما يتصل بها من الإيمان والعمل والوفاء والقصاص.

\separator

\section{تأويل (يا أيها الذين آمنوا كلوا من طيبات ما رزقناكم)}
\isnad{قال أبو جعفر:} هذا أمر من الله للمؤمنين أن يأكلوا من حلال الرزق وطيبه، وأن يدعوا ما كان أهل الجاهلية يحرمونه على أنفسهم بغير سلطان من الله، كالبحائر والسوائب ونحوها.

وبيّن \rawi{الطبري} أن الطيب هنا هو ما أحله الله لعباده فصار طيباً لهم بتحليله، لا بمجرد استحسان الناس أو عوائدهم.

\begin{fawaid}[الشكر لا يختزل في الثناء]
فسر الطبري الشكر هنا بالثناء على الله، وهذا من معاني الشكر لا تمامه. والشكر الشرعي أوسع من ذلك، إذ يدخل فيه استعمال النعمة في طاعة الله، كما قال تعالى: \textit{اعْمَلُوا آلَ دَاوُودَ شُكْرًا}. فالشكر قول وعمل، لا مجرد حمد باللسان.
\end{fawaid}

\separator

\section{تأويل (إنما حرم عليكم الميتة والدم ولحم الخنزير وما أهل به لغير الله)}
\isnad{قال أبو جعفر:} لم يحرم الله على المؤمنين من المطاعم إلا ما نص عليه في هذه الآية، فلا يجوز لهم أن يحرموا من عند أنفسهم ما لم يحرمه الله.

وتوسع \rawi{الطبري} في البيان اللغوي والإعرابي للآية، وذكر بعض الأوجه القرائية والعربية التي لها وجه في العربية، لكنه لا يجيز القراءة بها ما دامت مخالفة لإجماع القراءة المنقولة.

\begin{fawaid}[ثبوت الوجه العربي لا يكفي لإثبات القراءة]
من منهج الطبري أنه قد يقرر أن لبعض الألفاظ أو الأوجه وجهاً صحيحاً في العربية، لكنه مع ذلك لا يجيز القراءة به إذا لم تثبت به الحجة من قراءة الأمة. فصحة العربية شيء، وثبوت القرآن المتعبد بتلاوته شيء آخر.
\end{fawaid}

\subsection{معنى (وما أهل به لغير الله)}
فسر \rawi{الطبري} الإهلال هنا بما ذُبح لغير الله، أو ما ذُكر عليه غير اسم الله، وذكر أن أصل الإهلال رفع الصوت، ومنه إهلال المحرم بالتلبية، واستهلال الصبي عند ولادته.

ونبه الشيخ إلى أن \rawi{الطبري} استثنى ذبائح اليهود والنصارى من دخولها في هذا الباب، لأن الآية في ذبائح أهل الأوثان لا في طعام أهل الكتاب الذي جاءت الشريعة بإباحته.

\separator

\section{تأويل (فمن اضطر غير باغ ولا عاد فلا إثم عليه)}
\isnad{قال أبو جعفر:} من اضطر إلى أكل شيء من هذه المحرمات، لا باغياً ولا متعدياً، فلا حرج عليه في الأكل بقدر ما يدفع به الضرورة.

وجمع \rawi{الطبري} أقوال السلف في معنى \textit{غير باغ ولا عاد}، وناقش ما يدخل في البغي والعدوان، ثم قرر أن المراد نفي الإثم عمن التزم حد الضرورة ولم يتجاوزها.

\begin{fawaid}[الاضطرار لا يرفع الحكم مطلقاً بل بقدره]
الرخصة عند الضرورة لا تعني انحلال القيد كله، بل تقدر بقدرها. فمن اضطر إلى المحرم أبيح له منه ما يدفع ضرورته، بشرط ألا يكون باغياً أو متعدياً.
\end{fawaid}

\separator

\section{تأويل (إن الذين يكتمون ما أنزل الله من الكتاب ويشترون به ثمناً قليلاً)}
\isnad{قال أبو جعفر:} المراد بهذه الآية أحبار اليهود الذين كتموا صفة محمد ﷺ ونبوته، وأخذوا على ذلك ثمناً قليلاً من دنيا الناس.

وربط الشيخ هذه الآية بما قبلها من آية كتمان العلم، ونبه إلى أن من دقائق التفسير هنا النظر في مجموع الآيات لا في الآية المفردة وحدها.

\begin{fawaid}[الجمع بين الآيات يعين على تمام المعنى]
قد يترك الطبري إعادة بعض المعاني إذا كان قد مضى تقريرها في نظيرها من الآيات، فيفهم تفسير الآية اللاحقة على ضوء السابقة. وهذا يربي طالب العلم على قراءة السياق العام لا على تقطيع القرآن إلى مسائل منفصلة.
\end{fawaid}

\subsection{تأويل (أولئك ما يأكلون في بطونهم إلا النار)}
فسر \rawi{الطبري} ذلك بأن ما أخذوه من الرشوة والثمن على كتمان الحق يصير ناراً عليهم يوم القيامة، وأن الله لا يكلمهم كلام رضى، ولا يزكيهم من ذنوبهم، ولهم عذاب أليم.

\subsection{تأويل (فما أصبرهم على النار)}
ذكر \rawi{الطبري} الأقوال في معنى الآية، كقول من جعلها بمعنى: ما أجرأهم على العمل الذي يوجب النار، أو أنها خرجت مخرج التعجب من حالهم. ثم أدار المعاني بما يتسق مع أسلوب العرب.

\separator

\section{تأويل (ذلك بأن الله نزل الكتاب بالحق)}
بيّن \rawi{الطبري} أن سبب هذا الوعيد العظيم هو أن الله أنزل الكتاب بالحق، فخالفوه وكتموه واختلفوا فيه، فاستحقوا الشقاق البعيد.

\separator

\section{تأويل (ليس البر أن تولوا وجوهكم قبل المشرق والمغرب)}
جمع \rawi{الطبري} الأقوال في هذه الآية الجامعة، وذكر أن البر ليس محصوراً في مجرد التوجه إلى جهة من الجهات، بل حقيقة البر اجتماع أصول الإيمان وأعمال الطاعة الظاهرة والباطنة.

وساق أقوال من قال إن المعنى: ليس البر الصلاة وحدها، أو ليس البر مجرد استقبال القبلة، ثم رجح ما يجمع ذلك ويقرر أن البر اسم جامع لخصال الإيمان والطاعة.

\begin{fawaid}[الإيمان لا يصدق إلا بعمل]
هذه الآية من أوضح ما يرد به على المرجئة؛ لأن الله لما ذكر خصال الإيمان والإنفاق والوفاء والصبر قال بعدها: \textit{أُولَٰئِكَ الَّذِينَ صَدَقُوا}. فدل على أن صدق الإيمان لا يكون بمجرد دعوى القلب أو اللسان، بل بتحققه في العمل.
\end{fawaid}

\subsection{تأويل (وآتى المال على حبه)}
ذكر \rawi{الطبري} أن من تمام البر بذل المال مع محبته وشدة الحاجة إليه، فيؤتيه صاحبه لذوي القربى واليتامى والمساكين وابن السبيل والسائلين وفي الرقاب.

ونقل الشيخ هنا معنى نفيساً: أن هذا الحق في المال أعم من الزكاة، وأن الآية تدل على أنواع من الواجبات المالية والحقوق الشرعية التي لا تنحصر في قدر الزكاة المفروضة وحده.

\begin{fawaid}[في المال حقوق تتجاوز صورة الزكاة المجردة]
آية البر تدل على أن صاحب المال قد يجب عليه في بعض الأحوال بذل المال في وجوه البر والحق، ولا يحسن حصر الواجب المالي في صورة الزكاة فقط دون النظر إلى سائر الحقوق التي دل عليها الشرع.
\end{fawaid}

\subsection{تأويل (والموفون بعهدهم إذا عاهدوا)}
فسر \rawi{الطبري} العهد هنا بالوفاء بما أوجبه الله على عباده وبما يعاهد بعضهم بعضاً عليه من الحقوق والذمم.

\subsection{تأويل (والصابرين في البأساء والضراء وحين البأس)}
بيّن أن البأساء هي شدة الفقر، والضراء ما يصيب الإنسان في نفسه من مرض وبلاء، وحين البأس وقت القتال والحرب.

\begin{fawaid}[النصب على المدح من أساليب العربية في الثناء]
من المواضع اللغوية التي يبرزها الطبري هنا نصب \textit{الصابرين} على المدح، وهو من أساليب العرب في تخصيص بعض الصفات بمزيد ثناء وتنويه، مع أنها داخلة في السياق العام للصفات قبلها.
\end{fawaid}

\subsection{تأويل (أولئك الذين صدقوا وأولئك هم المتقون)}
\isnad{قال أبو جعفر:} من اتصف بهذه الخصال فهو الذي صدق الله في إيمانه، وهو الذي اتقى عقابه، لا من ادعى ذلك مع مخالفة أمره.

ونقل الشيخ عن \rawi{الحسن} و\rawi{الربيع} معنى عظيماً: أن حقيقة الإيمان العمل، وأن القول إذا لم يصحبه عمل لم يكن له حقيقة.

\separator

\section{تأويل (يا أيها الذين آمنوا كتب عليكم القصاص في القتلى)}
\isnad{قال أبو جعفر:} معنى الآية أن الله فرض عليكم القصاص في القتلى، أي: ألا تتعدوا بالقتل إلى غير القاتل، فيقتل الحر القاتل بالحر الذي قتله، والعبد القاتل بالعبد، والأنثى القاتلة بالأنثى، من غير مجاوزة ولا بغي.

وذكر \rawi{الطبري} في تفسير الآية أقوالاً متعددة في سبب نزولها وفي معنى ألفاظها:
\begin{itemize}
    \item قيل: نزلت في قبائل من العرب كانت تبغي في الجاهلية، فإذا قتل عبد منهم عبداً لم يرضوا حتى يقتلوا به حراً، وإذا قتلت امرأة امرأة لم يرضوا حتى يقتلوا بها رجلاً.
    \item وقيل: نزلت في طائفتين وقع بينهما قتال، فأمر النبي ﷺ أن تقع المقاصة بين ديات قتلاهم.
    \item وقيل: فيها معنى المساواة في القصاص مع البحث في الفضل بين دية الحر والعبد، والذكر والأنثى.
\end{itemize}

ثم ناقش \rawi{الطبري} هذه الأقوال ورد ما لم يقم عليه دليل صحيح أو خالف ما ثبت بالنص والإجماع، ورجح أن الآية ليست لنفي قتل الحر بالعبد أو الذكر بالأنثى، بل المقصود بها منع مجاوزة الجاني إلى غيره، ورفع بغي الجاهلية في التعدي بالقصاص.

\begin{fawaid}[سبب النزول من أول ما يطلب في التفسير]
من أهم ما يرفع الإشكال عن الآية معرفة سبب نزولها. فالطبري كثيراً ما يورد الاعتراض على ظاهر الآية، ثم يدفعه بالنظر في السبب والسياق، فيظهر بذلك وجه الكلام وينتظم المعنى.
\end{fawaid}

\begin{fawaid}[الترجيح بالإجماع والنص القاطع]
رد الطبري بعض الأقوال في آية القصاص بالاعتماد على الأخبار المتظاهرة عن النبي ﷺ، وعلى إجماع الأمة، فدل ذلك على أن الأقوال التفسيرية إذا خالفت النص القاطع أو الإجماع لم يجز حمل الآية عليها ولو كان لها وجه محتمل.
\end{fawaid}

\begin{fawaid}[القرآن يفسر بعضه بعضاً]
من محاسن هذا الموضع أن الطبري ومن قبله من السلف لم يقفوا عند آية القصاص وحدها، بل ربطوا بينها وبين قوله تعالى: \textit{وَكَتَبْنَا عَلَيْهِمْ فِيهَا أَنَّ النَّفْسَ بِالنَّفْسِ}. فاجتماع الآيتين يرفع الإشكال، ويبين أن المقصود منع بغي الجاهلية لا نفي القصاص بين أصناف الأحرار والعبيد والرجال والنساء على الإطلاق.
\end{fawaid}

\subsection{معنى (كتب عليكم)}
ذكر \rawi{الطبري} أن \textit{كتب} هنا بمعنى: فُرض وأُلزم، وهو استعمال معروف في لسان العرب وفي الشعر، ثم ربط ذلك بأصل الكتابة في اللوح المحفوظ، فكأن ما فرضه الله على عباده قد كُتب عليهم في سابق علمه وكتابه.

\subsection{معنى (القصاص)}
فسر \rawi{الطبري} القصاص بأنه المماثلة في الفعل، فقتل القاتل بمن قتل به قصاص؛ لأنه فُعل به مثل ما فعل، مع كون أحد الفعلين عدواناً والآخر عدلاً بحق.

\subsection{تأويل (فمن عفي له من أخيه شيء)}
ابتدأ \rawi{الطبري} هنا بذكر الأقوال في معنى العفو:
\begin{itemize}
    \item فقيل: المراد من عفا ولي الدم عن القصاص إلى الدية، فليكن طلبه لها بالمعروف، وعلى القاتل أداؤها بإحسان.
    \item وقيل: المراد ما فضل وبقي من الدية أو الأرش بعد المقاصة بين الحقوق.
\end{itemize}

ثم رجح أن الآية في العفو عن القصاص إلى الدية، وأن معنى الكلام: فمن صفح له ولي الدم عن القود إلى دية يأخذها منه، فليكن الاتباع بالمعروف والأداء بالإحسان.

\begin{fawaid}[تحرير موضع الخلاف في ألفاظ الآية]
يظهر في هذا الموضع تميز الطبري في تفكيك الخلاف: فهو لا يكتفي بقول ``اختلفوا''، بل يبين أن الخلاف قد يقع في سبب النزول، أو في معنى \textit{كتب عليكم}، أو في مدلول \textit{الحر بالحر}، أو في تفسير \textit{فمن عفي له}. وبهذا يحرر النزاع الحقيقي في كل جزء من الآية.
\end{fawaid}

\separator

وقف الشيخ عند هذا الموضع من آيات القصاص، وجعل على الطلاب واجباً في تلخيص الأقوال الواردة في سبب النزول وفي أجزاء الآية، وبيان الوجه الذي رجحه \rawi{الطبري} وسبب ترجيحه.
